
\subsection{Methodological Approach and Structure}
The review follows structured IS/medical literature guidance (concept
matrix inspired by \parencite{websterAnalyzingPrepareFuture2002} and PRISMA
principles). The methodological approach uses a multi‑stage search and
transparent extraction/synthesis procedures as outlined below.

\begin{enumerate}
  \item \textbf{Search strategy (Multi-stage, 2015--2025):}
    \begin{itemize}
      \item \textit{Scoping:} Google Scholar for broad coverage and
        identifying candidate papers.
      \item \textit{Formal queries:} IEEE Xplore, PubMed, and Scopus
        using controlled keyword strings and deduplication.
      \item \textit{Citation chasing:} Backward and forward citation
        tracking on key papers to ensure completeness.
    \end{itemize}
  \item \textbf{Inclusion / Exclusion criteria:}
    \begin{itemize}
      \item \textit{Inclusion:} Studies using wearable/non-EEG autonomic
        signals for seizure detection or preictal prediction.
      \item \textit{Exclusion:} EEG-only studies, invasive recordings,
        or papers limited to purely algorithmic simulations without
        human data.
    \end{itemize}
\end{enumerate}


Extracted dimensions populate a concept matrix with the following
fields:
\begin{enumerate}
  \item Setting: clinical vs ambulatory.
  \item Modality: ECG/HRV, PPG, EDA, ACC.
  \item Task: detection vs prediction.
  \item Model class and fusion strategy.
  \item Study phase: retrospective, pseudo‑prospective, prospective.
  \item Evaluation granularity: sample‑ vs alarm‑based.
  \item Windowing specifics (window length, stride, overlap).
  \item Metrics: sensitivity, FAR/h, AUC, f1-score, latency.
  \item Personalization.
  \item Interpretability and safety considerations.
  \item Data characteristics and dataset identity.
  \item Missing data handling (imputation, exclusion) and synthetic data use.
  \item Inference requirements (latency, computational load, hardware).
\end{enumerate}

\noindent\textbf{Prioritization Criteria:} 
Studies including splits preserving patient-specific data, FAR/h reporting 
and external or cross validation will be prioritized in the synthesis. A PRISMA-style flow diagram
will summarize and visualize screening and inclusion.

\subsection{Provisional Table of Contents}
\textit{Target Length: Approx. 15 Pages of Text}

\begin{enumerate}
  \item \textbf{Introduction} (approx. 1.5 pages)
    \begin{itemize}
      \item 1.1 Problem Statement and Motivation
      \item 1.2 Research Aim and Questions
      \item 1.3 Structure of this Review
    \end{itemize}
    
  \item \textbf{Methodological Approach} (approx. 1.5 pages)
    \begin{itemize}
      \item 2.1 Search Strategy (Multi-stage, 2015--2025)
      \item 2.2 Inclusion and Exclusion Criteria
      \item 2.3 Extraction Dimensions and Concept Matrix
      \item 2.4 PRISMA Flow Diagram
    \end{itemize}
    
  \item \textbf{Theoretical Background and Clinical Concepts} (approx. 2 pages)
    \begin{itemize}
      \item 3.1 Clinical Deployment Requirements (Sensitivity vs. FAR/h, Alarm Fatigue)
      \item 3.2 Seizure Classification and Clinical Implications
      \item 3.3 Autonomic Markers and Wearable Sensing (ECG/HRV, PPG, EDA, ACC)
      \item 3.4 Study Phases: Retrospective vs. Pseudo-Prospective vs. Prospective
    \end{itemize}
    
  \item \textbf{Modeling Landscape and Design Dimensions} (approx. 5 pages)
    \begin{itemize}
      \item 4.1 Architectural Evolution: Feature Engineering vs. Deep Learning
        \begin{itemize}
          \item Traditional Pipelines (SVM, RF, kNN with HRV/EDA Features)
          \item Deep Temporal Models (1D-CNN, LSTM, TCN)
          \item Attention and Transformer Variants
        \end{itemize}
      \item 4.2 Sensor Modalities and Fusion Strategies
        \begin{itemize}
          \item Feature-Level, Decision-Level, and Representation-Level Fusion
          \item Complementary Cues under Noise and Non-Stationarity
        \end{itemize}
      \item 4.3 Personalization and Domain Shift
        \begin{itemize}
          \item Patient-Specific Fine-Tuning and Calibration
          \item Cross-Dataset Evaluation and External Validation
        \end{itemize}
      \item 4.4 Robustness and Generalization
        \begin{itemize}
          \item Artifact Handling and Augmentation Strategies
          \item Out-of-Distribution Detection
        \end{itemize}
      \item 4.5 Deployment Constraints (Latency, Energy, Interpretability)
    \end{itemize}
    
  \item \textbf{Results: Synthesis and Comparative Analysis} (approx. 3 pages)
    \begin{itemize}
      \item 5.1 Performance Ranges by Study Phase and Setting
      \item 5.2 Model Comparison: Feature-Based vs. Deep Learning
      \item 5.3 Impact of Temporal Choices (Window Length, Stride, Context)
      \item 5.4 Evidence on Robustness and Domain Shift
      \item 5.5 Dataset Gaps, Class Imbalance, and Augmentation Effects
    \end{itemize}
    
  \item \textbf{Discussion} (approx. 1.5 pages)
    \begin{itemize}
      \item 6.1 Trade-offs: Complexity vs. Latency vs. Interpretability
      \item 6.2 Clinical vs. Ambulatory Deployment Considerations
      \item 6.3 Limitations of Current Evidence
    \end{itemize}
    
  \item \textbf{Conclusion and Future Work} (approx. 0.5 pages)
    \begin{itemize}
      \item 7.1 Summary of Key Findings
      \item 7.2 Practical Recommendations
      \item 7.3 Research Gaps and Future Directions
    \end{itemize}
\end{enumerate}

\vspace{1em}
\noindent\textbf{Note:} The following sections provide preliminary background 
to contextualize the research questions. They will be expanded with 
systematic literature findings in the final paper.
\vspace{0.5em}
